%  LaTeX support: latex@mdpi.com 
%  For support, please attach all files needed for compiling as well as the log file, and specify your operating system, LaTeX version, and LaTeX editor.

%=================================================================
\documentclass[water,article,submit,pdftex,moreauthors]{Definitions/mdpi} 

%=================================================================

% MDPI internal commands
\firstpage{1} 
\makeatletter 
\setcounter{page}{\@firstpage} 
\makeatother
\pubvolume{1}
\issuenum{1}
\articlenumber{0}
\pubyear{2024}
\copyrightyear{2024}
\externaleditor{Academic Editor: Firstname Lastname}
\datereceived{7 March 2024} 
\daterevised{27 March 2024} 
\dateaccepted{} 
\datepublished{} 
\hreflink{https://doi.org/} % If needed use \linebreak
%\doinum{}

%=================================================================
% Add packages and commands here. The following packages are loaded in our class file: fontenc, inputenc, calc, indentfirst, fancyhdr, graphicx, epstopdf, lastpage, ifthen, lineno, float, amsmath, setspace, enumitem, mathpazo, booktabs, titlesec, etoolbox, tabto, xcolor, soul, multirow, microtype, tikz, totcount, changepage, attrib, upgreek, cleveref, amsthm, hyphenat, natbib, hyperref, footmisc, url, geometry, newfloat, caption

\graphicspath{{figures/}} % path to folder with figures
\usepackage{subcaption}
\usepackage{listings}
\usepackage{xcolor}
\usepackage{gensymb} % for \textdegree
\usepackage{tabularx}
\usepackage{amsmath,mathtools,amsthm}
	\setcounter{MaxMatrixCols}{15}
\usepackage{array}
\usepackage{amssymb}
\usepackage{amsfonts}
\usepackage{float}
\usepackage{chngcntr}
\counterwithin*{equation}{section}
\usepackage[super]{nth}
\usepackage{longtable}
\usepackage{tabularx}
\usepackage{multirow}
\usepackage{adjustbox}

\definecolor{deepblue}{rgb}{0,0,0.5}
\definecolor{deepred}{rgb}{0.6,0,0}
\definecolor{deepmagenta}{RGB}{195,12,214}
\definecolor{deepgreen}{RGB}{29,122,30}
\definecolor{mintgreen}{RGB}{192,249,192}
\definecolor{codepurple}{RGB}{204, 0, 153}
\definecolor{LightCyan}{rgb}{0.88,1,1}
\definecolor{GainsboroPurple}{RGB}{247,176,247}
\definecolor{LightSlateGrey}{RGB}{124,118,118}
%    
\lstdefinestyle{mystyle}{
    commentstyle=\color{LightSlateGrey},
    keywordstyle=\color{deepgreen}\bfseries,
    emphstyle=\color{deepred},
    numberstyle=\tiny\color{deepblue},
    stringstyle=\color{codepurple},
    basicstyle=\ttfamily\scriptsize,
    breakatwhitespace=false,         
    breaklines=true,                 
    captionpos=b,                    
    keepspaces=true,                 
    numbers=left,                    
    numbersep=5pt,                  
    showspaces=false,                
    showstringspaces=false,
    showtabs=false,                  
    tabsize=2
}
\lstset{style=mystyle}

%=================================================================
% Full title of the paper (Capitalized)
\Title{Deep Learning Methods of Satellite Image Processing for Monitoring Flood Dynamics in the Ganges Delta, Bangladesh}

% MDPI internal command: Title for citation in the left column
\TitleCitation{Deep Learning Methods of Satellite Image Processing for Monitoring Flood Dynamics in the Ganges Delta, Bangladesh}

% Author Orchid ID: enter ID or remove command
\newcommand{\orcidauthorA}{0000-0002-5759-1089} % Polina Add \orcidA{} behind the author's name

% Authors, for the paper (add full first names)
\Author{Polina Lemenkova $^{1}$*\orcidA{}}

%\longauthorlist{yes}

% MDPI internal command: Authors, for metadata in PDF
\AuthorNames{Polina Lemenkova}

% MDPI internal command: Authors, for citation in the left column
\AuthorCitation{Lemenkova, P.}
% If this is a Chicago style journal: Lastname, Firstname, Firstname Lastname, and Firstname Lastname.

% Affiliations / Addresses (Add [1] after \address if there is only one affiliation.)
\address{%
$^{1}$ \quad Department of Geoinformatics, Faculty of Digital and Analytical Sciences, Universität Salzburg, Schillerstraße 30, A-5020 Salzburg, Austria.
}

% Contact information of the corresponding author
\corres{Correspondence: polina.lemenkova@plus.ac.at; Tel.: +43-677-6173-2772 (P.L.)}

% Current address and/or shared authorship
%\firstnote{Current address: Affiliation 3.} 

% Abstract (Do not insert blank lines, i.e. \\) 
\abstract{Mapping spatial data is essential for monitoring flooded areas, prognosis of hazards and prevention of flood risks. Ganges Delta\textcolor{blue}{, Bangladesh,} is the world's largest river delta prone to floods which impact social-natural systems through losses of lives, damage to infrastructure and landscapes. Millions of people living in this region are vulnerable \textcolor{blue}{to repetitive floods} due to exposure, high susceptibility and low resilience. Cumulative effects from monsoon climate, repetitive rainfalls and tropical cyclones, and hydrogeologic setting in Ganges Delta increase probability of floods. 
\textcolor{blue}{While engineering methods of flood mitigation include practical solutions (technical construction of dams, bridges, hydraulic drains), regulating traffic, land planning support systems, the geoinformation methods relies on modelling remote sensing (RS) data for evaluating the dynamics of flood hazards. Geo-information is indispensable for mapping catchment of flooded areas and visualization of affected regions in real time flood monitoring, implementing and developing emergency plans and vulnerability assessment through warning systems supported by RS data. With this regards, }this study \textcolor{blue}{ uses RS data to monitor} southern segment of Ganges Delta. \textcolor{blue}{The multispectral satellite }Landsat 8-9 OLI/TIRS satellite images \textcolor{blue}{were evaluated} in flood (March) and post-flood (November) periods for analysis of flood extent and landscape changes. \textcolor{blue}{Deep Learning} (DL) algorithms of GRASS GIS and modules of qualitative and quantitative analysis \textcolor{blue}{were used as advanced methods of satellite image processing}. The results presented a series of maps \textcolor{blue}{based} on the \textcolor{blue}{classified} images for monitoring floods \textcolor{blue}{in Ganges Delta}.}

% Keywords
\keyword{machine learning; deep learning; flood; monsoon; cartography; geoinformatics; remote sensing; satellite image; geospatial analysis; GRASS GIS} 

\PACS{91.10.Da; 91.10.Jf; 91.10.Sp; 91.10.Xa; 96.25.Vt; 91.10.Fc; 95.40.+s; 95.75.Qr; 95.75.Rs}
\MSC{86A30; 86-08; 86A99; 86A04}
\JEL{Y91; Q20; Q24; Q23; Q3; Q01; R11; O44; O13; Q5; Q51; Q55; N57; C6; C61}

\begin{document}

%----------- section ----------->
\section{Introduction}

\subsection{Background}
\hl{In hydrological studies and coastal management, detecting areas prone to floods is essential issue. Cartographic visualization of flooded landscapes using Remote Sensing (RS) data aims at flood hazard assessment and includes the analysis of the degree of flood impact, evaluation of the affected area and estimation of changes in the post-flood periods} \cite{su142114145,w15152802,su142316270}. \hl{Such analysis enables to highlight the spatial extent of the flood and spatio-temporal dynamics and has been acknowledged as one of the most powerful approaches in hydrological hazard monitoring. Indeed, satellite image analysis enables to reveal information about the consequences of flood as environmental changes such as affected land patches, destroyed areas, declined vegetation coverage and disrupted ecosystems} \cite{ijgi12110465,rs15102561}. 

\hl{Specifically for coastal Asian countries affected by the monsoon seasonality, fieldwork to capture data on flooded areas is difficult from both financially and for practical reasons, and, depending on the safety of the inundated area and severity of hydrological hazards, might become impossible. Therefore, RS data used for mapping flooded landscapes and spatial dynamics of water extent are critical tools for monitoring environmental consequences of hydrological hazards} \cite{YASEEN2024101665,GHALEHTEIMOURI2023}. 

\subsection{Related works}

Evaluating \hl{the }degree and rates of floods \hl{for hazard} prediction \hl{and mitigation of hydrological risks }can be performed using\hl{ }mage processing. \hl{Satellite images used for} mapping\hl{ }pre- and post-flood periods detected \hl{are valuable sources of information which can be retrieved using methods of GIS analysis and machine learning (ML)}. ML emerge as a fundamental technique of image processing and classification\hl{ }in a wide variety of applications including\hl{ }environmental studies \cite{ROCHA2022107415}, Earth science \cite{LI2023103345,Lemenkova202306,XIONG2021145256}, geological and hazard risk assessment \cite{AOURAGH2023100939}\hl{, hydrological monitoring} \cite{RAFIK2023101569} \hl{or} vegetation analysis \cite{WU2023110723,Lemenkova202310,DOSSANTOS2022106753}. \hl{Technical advantages of} ML \hl{have} been noted earlier \cite{9883389,coasts4010008}. These include\hl{ }high sensitivity of ML which enables to detect variations in the satellite images to evaluate environmental processes, and flexibility of image analysis which can be finely adjusted using diverse techniques of ML and algorithms. For example, existing ML \hl{algorithms} include Random Forest (RF), Support Vector Machines (SVM), etc.

\hl{Processing RS data to retreave information on floods from the satellite images is possible at diverse Geographic Information System (GIS). Diverse tools included in GIS allow mapping the degree and complexity of flood hazards and environmental consequences using classification of the satellite images. For instance, for coastal regions in deltaic areas prone to floods, time series of RS data can be processed for visualising dynamics of the ecological patterns related to flood events and monsoon cycles} \citep{RAZAVITERMEH2023100595,SINGHA2023316,MATHEW2021100861}. \hl{Moreover, the analysis of satellite images covering countries located close to shoreline can help to better understand the effects from climate changes on coastal landscapes and climate-environmental patterns of floods as well as to evaluate spatio-temporal dynamics in pre- and post-flood periods. To mention a few cases, satellite images have long been used for hydrological applications and geo-information, such as numerical evaluation of risks} \citep{THAKUR2024130683}, \hl{hydrological and irrigation modelling} \cite{SUN2024415}, \hl{susceptibility mapping of the coastal areas} \cite{RAZAVITERMEH2023100595}, \hl{detecting affected areas and estimating social-economic costs of floods} \cite{ZHENG2023101410,DUAN2024131010}, \hl{hydrogeological modelling} \cite{BELL2021107451}, \hl{land use mapping and evaluating climate and environmental impacts on vegetation} \cite{JARAMILLO2024526} \hl{or computing inundated areas} \cite{ZUO2024}.

\subsection{Gap and motivation}

\hl{Flood mitigation measures include diverse approaches and methods. Engineering solutions include, for instance, construction of infrastructure, dams, and bridges, maintaining the hydraulic performance of drains. Regulating traffic based on geo-information regarding location of buildings and land planning support systems enables to support affected communities. The non-engineering methods mostly include modelling, prediction and prognosis methods for implementing and developing emergency plans. Those enable to assess the vulnerability of area to flood hazards through warning systems supported by up-to-date geo-information. With this regards, using remote sensing (RS) data processed by GIS is indispensable for evaluating catchment flood areas and visualization of the effected regions in real time flood monitoring.}

Currently, the cartographic monitoring of flood in Ganges Delta presents efforts on GIS-based mapping\hl{ }\cite{Sanyal2023,ISLAM2023167,RANA2023}. \hl{Besides the conventional approaches, recent studies in the same area study area use alternative methods such as Sentinel-1 SAR images ad Google Earth Engine} \cite{ChakmaAkter,SINGHA2020278}. \hl{Existing approaches} raised the question of\hl{ the }operative and accurate \hl{methods for }data processing \hl{aimed at} flood prediction and management\hl{. }Nevertheless, operative monitoring of floods require advanced machine learning techniques and programming approach to support real-time flood prognosis systems through the algorithms of artificial intelligence \cite{ZHAO2021126777}. Identifying landscape patterns within coastal and tidal areas, estuary and delta is challenging due to the obscurity and ambiguity of land cover types. Accurate classification of the satellite images for environmental mapping requires advanced methods of scripting languages which aim at robust detecting of the patterns \cite{YANG2020100413,Lemenkova202322,BEVERIDGE2020104843,SADIQ2023111233,Lemenkova202227}. Existing state-of-the-art methods \hl{of image processing and classification }mostly use the traditional GIS for \hl{data handling} which may lead to misclassification and inaccurate labelling of pixels. 

\hl{Among the available traditional tools of satellite image processing and classification, novel methods of Deep Learning (DL) can now be used as a complementary advanced approach in RS. For instance, such algorithms as Multilayer Perceptron (MLPClassifier) are designed to support image partition, classification and analysis through algorithms of decision trees and computer vision} \cite{6493388,6422192,7507955}. \hl{The application of such methods in hydrological studies and coastal risk assessment creates principally new perspectives in cartography by combining RS with programming approaches for mapping flooded areas. It is therefore particularly important to apply such advanced technical tools to regions with high heterogeneity and complexity of landscape patterns, such as coastal zones, where algorithms of DL can generate decisions on image classification and detect affected areas}. 

When applied to a series of satellite images, ML methods can effectively discover the properties of landscapes through comparative analysis \cite{9324433,Lemenkova202313}. Moreover, ML enables to quantify the patches of landscapes \hl{on} raster images \hl{using} algorithms \hl{of} computer vision  \cite{10281789,jmse11040871,8899013,Lemenkova202021,9333522} \hl{or} determine landscape dynamics for environmental monitoring \cite{9392684,Lemenkova202324}. Finally, other advantages of the ML approach is that it is a resource- and time-effective method based on advanced programming methods which largely involves scripting techniques. Such approach enables to automate data processing and cartographic analysis as mentioned earlier \cite{10282693,Lemenkova202301,10441504}. Hence, using ML enables to indicate flood risk and\hl{ }development of crisis situation in flood-prone areas such as delta of Ganges, Bangladesh. Moreover, such methods are especially effective for detecting correlation between \hl{the }environmental processes and climate change over time which is especially actual for tropical regions that are subject to monsoon. 

\subsection{Goals and objectives}

\hl{This paper presents the use of the advanced deep learning (DL) tools of GRASS GIS for image classification with a case of Ganges River Delta. Specifically, several Landsat satellite images are used, processed, classified, compared and analysed to evaluate pre-and post flooded periods in coastal Bangladesh. The potential of the ML modules of GRASS GIS for mapping consequences of hydrological hazards is exemplified with the case of affected landscapes of coastal Bangladesh and associated climate dynamics related to the monsoonal climate of the Indian Ocean.} \hl{The focus of the manuscript is on a holistic approach to accurate satellite image processing with the help of the open source software GRASS GIS. The methodology applied to the subject of the study can be applied to the other relevant countries located in coastal regions with flood-prone areas.}

\section{\textcolor{blue}{Study Area}}

\subsection{Current landscape}

This study presents a case of the Ganges River Delta, located in southern Bangladesh \hl{and prone to floods}, Figure \ref{fig01}. The Ganges Delta is the world's largest river delta with hundreds of millions people living around \hl{in a total catchment area}. Located in the region with monsoon climate and seasonally changing environmental setting, it is subject to regular floods, occasional tropical cyclones, tidal waves and related inundations \cite{Bhardwaj2022}. Complex hydrology of the Ganges Delta increases the consequences of floods through dense network of streams, rivers and tributaries where water remains for a longer period \cite{Islam}. Alluvial clayey sediments dominating in the riverbed\hl{ }contribute to the stagnation of water during flood \hl{because they } accumulate in \hl{numerous interconnected} channels, inner lakes, wetlands and flood plains of the Ganges River Delta \cite{Datta,KHAN201927,Lemenkova202029}. The monsoon climate regulates the repeatability of rainfalls and floods in Ganges Delta which normally augment between the months of March and September. 

\begin{figure}[H]
    \begin{center}
	\hspace{-35pt}\includegraphics[width=11.0cm]{Fig_01.jpg}
    \end{center}
    \caption{Topographic map of Bangladesh with depicted study area showing the placement of the Landsat data over the Ganges River Delta. \hl{Digital elevation data: SRTM/GEBCO, 15 arc sec resolution grid.} Software: GMT version 6.4.0. Map source: author.}
    \label{fig01}
\end{figure}

\subsection{Problem formulation}
The repetitive floods in Ganges Delta negatively impact social-natural systems causing losses of lives and damage to infrastructure \cite{JERIN2023}, fragmentation and disruption in natural landscapes \cite{MUKHERJEE2021106961} \hl{and} salinity intrusion during storms\hl{. Besides, floods affect} crop agriculture \cite{ABEDIN2023315,ALMASUD2020104443}, and \hl{negatively influence the} sustainability of coastal ecosystems \cite{SARKER2024169718,MAINUDDIN2021105740,HASAN2023101028}. Millions of people living in southern Bangladesh and neighbouring regions of India are vulnerable \hl{to repetitive floods} due to exposure, high susceptibility and low resilience \cite{SINGHA2020106825,RUMPA2023104047}. 

The examples of social effects from floods include the disruptions to clean drinking water and electricity, broken transport and communication systems, effects regular education and normal working environment and health care support for population \cite{JERIN2023103851,NAHIN2023e14520}. The long-term effects from flood and its consequences which may last for months  seriously affect the social-economic system of Bangladesh. The coastal area of Ganges River Delta supports approximately 56.7 M of \hl{population} \cite{DAS2021101983} where many \hl{settlements} are \hl{located including both small villages and large cities such as the capital Dhaka, Khulna (the third-largest city in the country), Comilla and Narayani. Supporting such flood-prone areas }requires\hl{ geospatial }monitoring of flood risks \hl{in order to visualise the affected areas on the maps of floods} \cite{Alietal,AZAD2023100498}.

The uncertainty of flood prediction in Ganges River Delta has an impact on vulnerability to risk and information among the local population. Hence, preventive prediction rises a question of using reliable data and advances methods for operative monitoring of floods and assessment of environmental sustainability \cite{NICHOLLS2016370}. \hl{The main approach that a country can take to assess flood risk is to map areas that are susceptible to flooding based on their geography and the comprehensive flood risk assessment. With this regards, flood maps should be prepared by the relevant authorities, and remote flood monitoring should be used to reduce the risk during weather events. Various techniques have been developed to produce flood maps. Among them, satellite-based flood maps tend to provide real-time, contiguous representations with varying spatial resolution over time and higher accuracy.}

Notably, no prior study has explored the deep learning (DL) methods available in Geographic Resources Analysis Support System (GRASS) Geographic Information System (GIS) through a comprehensive evaluation of multi-spectral imagery for mapping floods in Ganges Delta\hl{, Bangladesh}. As a response to this gap, this study presents \hl{a report with} such approach which is based on using \hl{the }DL methods of \hl{satellite }image processing by GRASS GIS scripting software. \hl{The }implementation of this approach is realised through the MLPClassifier algorithm\hl{ }from the embedded Scikit-Learn library of Python. 

%----------- section ----------->
\section{Materials and Methods}

%----------- subsection ----------->
\subsection{Deep Learning}

As a branch of Machine learning (ML), Deep learning (DL) is increasingly used in Earth science for automatic processing of geospatial data. The problem remains with the technical approach of such methods, especially in Remote Sensing (RS) data processing. DL approaches can be applied by a variety of tools, using programming libraries, embedded plugins and modules in Geographic Information Systems (GIS) software. Amongst geoinformation software, GRASS GIS \cite{GRASS_GIS_software,NETELER2012124,MitasovaNeteler} proposes effective solutions to satellite image processing with the option to utilise diverse choices of models for classification of RS data, Figure \ref{fig02}.

\begin{figure}[H]
    \begin{center}
%    \setlength{\unitlength}{0.5cm}
	\hspace{-35pt}\includegraphics[width=8.0cm]{Fig_02.jpg}
    \end{center}
    \caption{Conceptual structure of the ML and DL algorithms in geospatial data analysis and image processing by GRASS GIS. Graphical software: Inkscape version 1.2 \cite{Inkscape}. Scheme source: author.}
    \label{fig02}
\end{figure}

Of the existing ML methods, the most widely used approach is the Deep Learning (DL) with many examples of applications for image processing in environmental studies \cite{10441152,AHMED2024167234,HASSAN2023100399}. DL presents advanced technique for evaluating the environmental dynamics during the climate changes assessed from image analysis \cite{10441576}. There are many categories of DL algorithms which have various setting parameters and functionality for measurements of pixels values in the image and classification. Overall, DL algorithms train the computer model using knowledge derived from the input layers and processed in hidden layers which results in accurate modelling as a result of data processing.

The existing algorithms use Neural Networks (NN) to learn semantic representations of features as indicated by pixels on the images hence derived directly from Earth Observation (EO) data. This enables to evaluate landscapes properties and dynamics of the environmental properties \cite{BISWAS2023e21245}, using such methods as segmentation \cite{Lemenkova202320}, feature extraction, and classification \cite{RAISA2024101128}. Further development of DL methods aims at achieving high precision of the image processing as signals obtained from the spectral reflectances of the land cover types detected on the space borne images, as well as the use of complex programming algorithms, for instance, such as in NN \cite{RUDRA2023e16459}. 

%----------- subsection ----------->

\subsection{Workflow}
The methodological workflow employed in this study consisted of several steps and processes which aimed at analysis of geospatial information and image processing. It also included diverse data obtained from different sources of information. While the purpose of Earth observation (EO) data is to provide information on objects and features visible from space, choice of methods serves varied purposes that accommodate specific DL algorithms and techniques of image processing \cite{4736879,5070195}. \hl{This study} \hl{utilises} the \hl{existing} DL \hl{methods available in the modules }of GRASS GIS for satellite image processing \hl{to detect} flooded areas in Ganges River Delta, southern Bangladesh. 

%----------- subsection ----------->
\subsection{Theoretical fundamentals of MLP}
The Multilayer perceptron is a class of DL and a conceptual theoretical background to data analysis performed in this study which should be briefly discussed. A multilayer perceptron (MLP) presents a feedforward Artificial Neural Network (ANN), a branch of DL methods \cite{su151813786}. It consists of fully connected neurons with a nonlinear activation function that is organised in at least three layers, Figure \ref{fig03}. 

\begin{figure}[H]
    \begin{center}
	\includegraphics[width=14.0cm]{Fig_03.jpg}
    \end{center}
    \caption{General methodological scheme for the Deep Learning (DL) approach used for satellite image classification and EO data processing. Software: R, library \hl{DiagrammeR} \cite{DiagrammeR}. Diagram source: author.}
    \label{fig03}
\end{figure}

These layers are being able to distinguish data that is not linearly separable and, for time series of satellite images, the variables that are varying in time. \hl{This results in different pixel values being recognised by the model as Digital Number (DN)}. Such characteristics is caused by the spatio-temporal variability on spectral reflectance of pixels that constitute the image as raster matrix according to the physical properties of landscape and vegetation patches \cite{rs16010184,HUANG20181915,YIN2021102514}. 

\subsection{Data}
In this study, we used \hl{eight} Landsat 8-9 Operational Land Imager (OLI) and Thermal Infrared Sensor (TIRS) multispectral satellite images for environmental mapping of Ganges Delta, Bangladesh. \hl{the images are selected} for flood and post-flood periods (March and November) for year 2021, 2022 and 2023. The \hl{images for March and November of} year \hl{2014} were used as \hl{a seed of} training pixels for Deep Learning classification. 

The generated dataset is \hl{collected} from the \href{https://earthexplorer.usgs.gov/}{EarthExplorer} \cite{USGSPublication} with image samples openly available for download, processing and data reuse. The original Landsat images included multispectral, panchromatic and SWIR bands, Figure \ref{fig04}. The images are in GeoTIFF format and contained the \hl{seven multispectral} bands for each image. Hence, the image recordings received from the USGS source include RS datasets have been explored for metadata using records available in \hl{the Appendix in }Table \ref{tabApp} \hl{and the essential characteristics are summarised in Table} \ref{tab01}. The advantage of the Landsat archives is that it is reusable. Thus, in future studies, these data can be used as input data for similar tasks, such as classification and detection of land cover types, image processing and segmentation using presented GRASS GIS workflow model of other tools. 

\begin{table}[H]
\footnotesize
\caption{\hl{Attribute table of the main characteristics of} the Landsat 8-9 OLI/TIRS images} \label{tab01}
\begin{adjustwidth}{-\extralength}{0cm}
    \begin{tabularx}{\fulllength}{llll}
    \toprule
       	\textbf{Data Set Attribute} & \textbf{Attribute Value} & \textbf{Attribute Value} & \textbf{Attribute Value} \\ \hline
        \textbf{Landsat Scene Identifier} & LC81370442023066LGN00 & LC81370442022079LGN00 & LC81370442021076LGN00 \\ \hline
        \textbf{Date Acquired} & \textbf{2023/03/07} & \textbf{2022/03/20} & \textbf{2021/03/17} \\ \hline
        \textbf{Land Cloud Cover} & 3.12 & 0.01 & 0.03 \\ \hline
        \textbf{Scene Cloud Cover L1} & 2.97 & 0.01 & 0.03 \\ \hline
        \textbf{Sun Elevation L0RA} & 51.70012312 & 56.10505202 & 55.19368850 \\ \hline
        \textbf{Sun Azimuth L0RA} & 134.86314148 & 129.90704446 & 131.01649112 \\ \hline
        \textbf{TIRS SSM Model} & FINAL & FINAL & FINAL \\ \hline
        \textbf{Data Type L2} & OLI\_TIRS\_L2SP & OLI\_TIRS\_L2SP & OLI\_TIRS\_L2SP \\ \hline
        \textbf{Satellite} & 8 & 8 & 8 \\ \hline\hline
        \midrule
        \textbf{Data Set Attribute} & \textbf{Attribute Value} & \textbf{Attribute Value} & \textbf{Attribute Value} \\ \hline
        \textbf{Landsat Scene Identifier} & LC91370442023330LGN00 & LC91370442022327LGN01 & LC81370442021332LGN00 \\ \hline
        \textbf{Date Acquired} & \textbf{2023/11/26} & \textbf{2022/11/23} & \textbf{2021/11/28} \\ \hline
        \textbf{Land Cloud Cover} & 0.02 & 0.20 & 0.40 \\ \hline
        \textbf{Scene Cloud Cover L1} & 0.02 & 0.19 & 0.38 \\ \hline
        \textbf{Sun Elevation L0RA} & 41.83496249 & 42.43660966 & 41.36291892 \\ \hline
        \textbf{Sun Azimuth L0RA} & 154.44654325 & 154.42300606 & 154.52745495 \\ \hline
        \textbf{TIRS SSM Model} & N/A & N/A & FINAL \\ \hline
        \textbf{Data Type L2} & OLI\_TIRS\_L2SP & OLI\_TIRS\_L2SP & OLI\_TIRS\_L2SP \\ \hline
        \textbf{Satellite} & 9 & 9 & 8 \\ \hline
        \bottomrule
    \end{tabularx}
    \end{adjustwidth}
\end{table}

The datasets are available in EarthExplorer \cite{USGSPublication} and originating from NOAA Landsat collections \cite{Landsat}. The images were pre-projected as original data and had the following technical characteristics: Datum and Ellipsoid World Geodetic System 84 (WGS84), stored in Universal Transverse Mercator (UTM) Zone 46 for southern Bangladesh, and Worldwide Reference System (WRS) Path/Row 137/44. The data were collected at Nadir in day time. The Landsat Collection Category T1, Number 2 for all the scenes. Landsat Station Identifier is LGN. Sensor Identifier is Landsat OLI TIRS for all the images, Data Type L2 is OLI TIRS L2SP. The remaining essential characteristics of the images are summarised in Table \ref{tab01}.

The selection of these data is explained by the reputation and reliability of the Landsat imagery widely used for environmental studies. Besides, the Landsat images are available free of charge, have high repeatability of coverage and high quality. Hence, satellite images present a reliable source of information for detecting the extent of floods and visualising inundated areas using time series analysis. Their applications is reported in many existing studies with a focus on environmental monitoring \cite{Sheonty,Islam2016,Lemenkova202229,Ferdous,Lemenkova202321}. The data were collected from the available freely from the United States Geological Survey (USGS) repository EarthExplorer \cite{USGSPublication}. The original data are visualised in Figure \ref{fig04}. 

\begin{figure}[H]
  \centering
  \begin{tabular}[b]{c}
    \includegraphics[width=2.9cm]{Fig_04a.jpg} \\
    \small (a): March 2014
  \end{tabular}
  \begin{tabular}[b]{c}
    \includegraphics[width=2.9cm]{Fig_04b.jpg} \\
    \small (b): November 2014
  \end{tabular}
    \begin{tabular}[b]{c}
    \includegraphics[width=2.9cm]{Fig_04c.jpg} \\
    \small (c): March 2021
  \end{tabular}
  \begin{tabular}[b]{c}
    \includegraphics[width=2.9cm]{Fig_04d.jpg} \\
    \small (d): November 2021
  \end{tabular}
   \begin{tabular}[b]{c}
    \includegraphics[width=2.9cm]{Fig_04e.jpg} \\
    \small (e): March 2022
  \end{tabular}
  \begin{tabular}[b]{c}
    \includegraphics[width=2.9cm]{Fig_04f.jpg} \\
    \small (f): November 2022
  \end{tabular}
    \begin{tabular}[b]{c}
    \includegraphics[width=2.9cm]{Fig_04g.jpg} \\
    \small (g): March 2023
  \end{tabular}
  \begin{tabular}[b]{c}
    \includegraphics[width=2.9cm]{Fig_04h.jpg} \\
    \small (h:) November 2023
  \end{tabular}
  \caption{Data sources: Landsat 8-9 OLI/TIRS images on Ganges Delta, Bangladesh, collected during the flood period (March) and post-flood period (November) on years 2014, 2021, 2022 and 2023, used for analysis of flood dynamics. Data source: USGS. Compilation source: author.}
  \label{fig04}
\end{figure}

\subsection{Software}
The workflow used in this research employs a chain of modules that use scripting GRASS GIS software \cite{GRASS_GIS_software} for image processing described in earlier works \cite{Lemenkova202323,technologies11020046}. For the installation and use of Machine Learning (ML) modules of GRASS GIS, we recommend using a Python version 3.8.5 or higher. Specifically, the installation of our framework can be performed using the available packages in the shared data. The additional Python libraries required by GRASS GIS are included in the modules through embedded links and include the following: Matplotlib, NumPy, Pandas, Scikit-Learn, SciPy. Additionally, the XCode is need to be installed for writing the scripts and using GitHub is recommended. These libraries can be installed using 'pip install' with more details in the online resources of the GRASS GIS. The methodology of time series analysis includes the unsupervised and supervised classification of multispectral satellite images for detecting changes in landscapes around the Ganges Delta in pre- and post-flood periods. The map in Figure \ref{fig01} is prepared using \hl{Generic Mapping Tools (GMT) version 6.4.0} \cite{Wessel,WesselSmith} using raster grid of General Bathymetric Chart of the Oceans (GEBCO). 

%----------- subsection ----------->
\subsection{Implementation}
We used a workflow in GRASS GIS platform as a container of diverse modules to run the scripts on MacOS. The GRASS GIS uses algorithms of ML/DL methods of image processing derived from the Python's Scikit-Learn library via the integrated libraries in the machine. For GRASS GIS platform, this is performed using the containers of modules via the sequence of commands presented below in scripts in Listings \ref{lst01}, \ref{lst02} and \ref{lst03}. The processes include the import and pre-processing of the multispectral bands of the given image, image analyse and classification using difference in spectral reflectance associated with land patches to each valid patch of the Ganges River Delta landscapes evaluated for March (flood) and November (post-flood) periods of 2021, 2022 and 2023.

The following is an outline of the essential details regarding the approaches and techniques of the algorithm that compose the GRASS GIS framework using for image analysis. For details, full codes are presented below and explained with essential details. In all the GRASS GIS modules and functions and post-processing of the satellite images using DL approach, it is necessary to apply the module which implement the required functionality for data processing for each image in terms of the identifiers of functions and parameters of image processing. For this purpose, we adopted and modified the approaches of ML modules (libraries) of GRASS GIS which rely on using the Python's Scikit-Learn library for DL techniques.

%----------- subsection ----------->
\subsection{Image processing}
The images were first imported into the GRASS GIS system using 'r.import' module and colour composites are created using 'r.composite'. The images were then displayed using the combination of modules 'd.mon' and 'd.rast'. The files are then saved to the bitmap format using the 'd.out.file' module. The implementation of this step is performed using the GRASS GIS syntax (example is presented for November 2023), presented in the codes of Listing \ref{lst01}.

\begin{lstlisting}[language=bash,caption=GRASS GIS code for data import and creating colour composites algorithm,style=mystyle,label={lst01}]
#importing the image subset with 7 Landsat bands and display the raster map
r.import input=/Users/polinalemenkova/grassdata/Bangladesh/LC09_L2SP_137044_20231126_20231128_02_T1_SR_B1.TIF output=L8_2023_N_01 extent=region resolution=region --overwrite
r.import input=/Users/polinalemenkova/grassdata/Bangladesh/LC09_L2SP_137044_20231126_20231128_02_T1_SR_B2.TIF output=L8_2023_N_02 extent=region resolution=region
r.import input=/Users/polinalemenkova/grassdata/Bangladesh/LC09_L2SP_137044_20231126_20231128_02_T1_SR_B3.TIF output=L8_2023_N_03 extent=region resolution=region
r.import input=/Users/polinalemenkova/grassdata/Bangladesh/LC09_L2SP_137044_20231126_20231128_02_T1_SR_B4.TIF output=L8_2023_N_04 extent=region resolution=region
r.import input=/Users/polinalemenkova/grassdata/Bangladesh/LC09_L2SP_137044_20231126_20231128_02_T1_SR_B5.TIF output=L8_2023_N_05 extent=region resolution=region
r.import input=/Users/polinalemenkova/grassdata/Bangladesh/LC09_L2SP_137044_20231126_20231128_02_T1_SR_B6.TIF output=L8_2023_N_06 extent=region resolution=region
r.import input=/Users/polinalemenkova/grassdata/Bangladesh/LC09_L2SP_137044_20231126_20231128_02_T1_SR_B7.TIF output=L8_2023_N_07 extent=region resolution=region
g.list rast
# false color
r.composite blue=L8_2023_N_07 green=L8_2023_N_05 red=L8_2023_N_03 output=L8_2023_N_753
d.mon wx0
d.rast L8_2023_N_753
d.out.file output=Bangladesh_753 format=jpg --overwrite
# false color: NIR band B05 in the red channel, red band B04 in the green channel and green band B03 in the blue channel
r.composite blue=L8_2023_N_03 green=L8_2023_N_04 red=L8_2023_N_05 output=L8_2023_N_345
d.mon wx0
d.rast L8_2023_N_345
d.out.file output=Bangladesh_345 format=jpg --overwrite
# true color
r.composite blue=L8_2023_N_02 green=L8_2023_N_03 red=L8_2023_N_04 output=L8_2023_N_234
d.mon wx0
d.rast L8_2023_N_234
d.out.file output=Bangladesh_J_234 format=jpg --overwrite
\end{lstlisting}

The results of the colour composites are demonstrated in Figure \ref{fig05}. 

\begin{figure}[H]
	\begin{subfigure}[b]{.45\textwidth}
		\centering
		\includegraphics[width=0.98\textwidth]{Fig_05a.jpg}
		\caption{Bands 3-4-7}
	\end{subfigure}%
	\begin{subfigure}[b]{.45\textwidth}
		\centering
		\includegraphics[width=0.98\textwidth]{Fig_05b.jpg}
		\caption{Bands 3-4-5}
	\end{subfigure}
		\hfill
		\vspace{1em}
	\begin{subfigure}[b]{.45\textwidth}
		\centering
		\includegraphics[width=0.98\textwidth]{Fig_05c.jpg}
		\caption{Bands 2-4-7}
	\end{subfigure}
	\begin{subfigure}[b]{.45\textwidth}
		\centering
		\includegraphics[width=0.98\textwidth]{Fig_05d.jpg}
		\caption{Bands 3-7-4}
	\end{subfigure}
\caption{Examples of different false colour composites generated using the multispectral bands of the Landsat 8-9 OLI/TIRS scenes by GRASS GIS 'r.composite' module. Here the band numbers correspond to the following channels of the Landsat 8-9 OLI/TIRS images: Band 2 -- Blue; Band 3 -- Green; Band 4 -- Red; Band 5 -- Near Infrared (NIR); Band 7 -- Shortwave Infrared-2 (SWIR-2).}\label{fig05}
\end{figure}

The unsupervised classification is based on clustering using k-means and Maximum Likelihood embedded in the GRASS GIS. technically, it was performed using the modules 'i.group', 'i.cluster ' and 'i.maxlik' which creates the signature file and performs accuracy assessment using 'reject' function that evaluates rejection probability of the pixels using chi-square test. The implementation of this step is performed using the GRASS GIS syntax (example is presented for November 2023), presented in the codes of Listing \ref{lst02}.

\begin{lstlisting}[language=bash,caption=GRASS GIS code for unsupervised image classification method using k-means clustering algorithm,style=mystyle,label={lst02}]
g.region raster=L8_2023_N_01 -p
i.group group=L8_2023_N subgroup=res_30m \
  input=L8_2023_N_01,L8_2023_N_02,L8_2023_N_03,L8_2023_N_04,L8_2023_N_05,L8_2023_N_06,L8_2023_N_07 --overwrite
# Clustering: generating signature file and report using k-means clustering algorithm
i.cluster group=L8_2023_N subgroup=res_30m \
  signaturefile=cluster_L8_2023_N \
  classes=10 reportfile=rep_clust_L8_2023_N.txt --overwrite
# Classification by i.maxlik module
i.maxlik group=L8_2023_N subgroup=res_30m \
  signaturefile=cluster_L8_2023_N \
  output=L8_2023_N_cluster_classes reject=L8_2023_N_cluster_reject --overwrite
\end{lstlisting}

The results of this step are demonstrated in Figure \ref{fig06} and are used as training data for supervised classification using deep learning approach. Moreover, the selected example of accuracy analysis for 2014 which was used as training data set is shown in Figure  \ref{fig07}. This step includes the procedures of clustering that generated signature file and report which was performed using k-means algorithm using 'i.cluster' module. The unsupervised classification was done using the 'i.maxlik' module. The results of this step are presented in Figure \ref{fig02} and provided the seed data and training map for the next step of DL modelling. 

\begin{figure}[H]
	\begin{subfigure}[b]{.50\textwidth}
		\centering
		\includegraphics[width=0.98\textwidth]{Fig_06a.jpg}
		\caption{March 2014}
	\end{subfigure}%
	\begin{subfigure}[b]{.50\textwidth}
		\centering
		\includegraphics[width=0.98\textwidth]{Fig_06b.jpg}
		\caption{November 2014}
	\end{subfigure}
\caption{Image processing of the Landsat 8 OLI/TIRS scene on March and November 2014 using unsupervised classification by clustering and Maximal Likelihood method: \textbf{(a)}: March 2014; \textbf{(b)}: November 2014.}\label{fig06}
\end{figure}

\begin{figure}[H]
	\begin{subfigure}[b]{.50\textwidth}
		\centering
		\includegraphics[width=0.98\textwidth]{Fig_07a.jpg}
		\caption{March 2014}
	\end{subfigure}%
	\begin{subfigure}[b]{.50\textwidth}
		\centering
		\includegraphics[width=0.98\textwidth]{Fig_07b.jpg}
		\caption{November 2014}
	\end{subfigure}
\caption{Accuracy assessment of image classification of the Landsat 8 OLI/TIRS scenes: \textbf{(a)}: March 2014; \textbf{(b)}: November 2014.}\label{fig07}
\end{figure}

Deep learning (DL) analysis is performed using algorithm of Multilayer Perceptron (MLPClassifier) for automated classification of the series of satellite images. First, the spatial region extent was defined using 'g.region' module which sets up the coordinates of the maps within the project using the target map. Afterwards, the training pixels from an older (2014) land cover classification were generated using 'r.random' module of GRASS GIS and used as seeds. Next, the imagery group with all Landsat-8 OLI/TIRS bands was created to include all multispectral bands (we do not need panchromatic bands in this case). Then the training pixels were applied to perform a classification using 'MLPClassifier' of deep learning approach on recent Landsat (in the presented code below -- the image on November 2023). This method trains a MLPClassifier model using 'r.learn.train. 

After this step, the prediction of the pixel assignments to each land cover class over the Ganges Delta River was performed using 'r.learn.predict' module. The raster categories were automatically applied to the classification output and checked again using the 'r.category' module. Then, the color scheme was assigned from the land class training map to result and visulised using the modules 'r.colors' and 'd.rast'. Technically, the DL approach was realised using the 'r.learn.train' module of GRASS GIS and demonstrated in script of Listing \ref{lst03}.

\begin{lstlisting}[language=bash,caption=GRASS GIS code for unsupervised image classification method using k-means clustering algorithm,style=mystyle,label={lst03}]
g.region raster=L8_2023_N_01 -p
r.random input=L8_2014_N_cluster_classes seed=100 npoints=1000 raster=training_pixels
i.group group=L8_2023_N input=L8_2023_N_01,L8_2023_N_02,L8_2023_N_03,L8_2023_N_04,L8_2023_N_05,L8_2023_N_06,L8_2023_N_07 --overwrite
r.learn.train group=L8_2023_N training_map=training_pixels \
    model_name=MLPClassifier n_estimators=500 save_model=mlpc_model.gz --overwrite
r.learn.predict group=L8_2023_N load_model=mlpc_model.gz output=mlpc_classification_Ganges --overwrite
r.category mlpc_classification_Ganges
d.mon wx0
d.rast mlpc_classification_Ganges
r.colors mlpc_classification_Ganges color=roygbiv -e
d.legend raster=mlpc_classification_Ganges title="MLPClassifier: 11/2023" title_fontsize=14 font="Helvetica" fontsize=12 bgcolor=white border_color=white
d.out.file output=MLPC_2023_11 format=jpg --overwrite
\end{lstlisting}

Each image of the produced dataset contained seven multispectral bands processed automatically by GRASS GIS modules 'r.learn.train' and 'r.learn.predict' for each image. The total time processing for one image is ca. 30 minutes, which demonstrates a high time efficiency of the proposed GRASS GIS algorithm of deep learning. The algorithm of MLPClassifier was tested initially on a single image processed as a complete workflow for datasets on March and November 2021, 2022 and 2023 and then analysed for changes in the evaluated period. Changes detected using the repeatability pattern of land cover types over the Ganges River Delta were evaluated for the flood- and post-flood periods. Noise background related to cloudiness is ignored since the images were collected with a high visibility (below 10\% of cloudiness) quality which ensured that land cover types are clearly visible and identified by the ML algorithm. 

%----------- section ----------->
%\pagebreak
\section{Results}

Satellite images Landsat evaluated in this study cover the regions of southern Bangladesh and cover periods of March and November for years 2021, 2022 and 2023. Land cover changes during flood- and post-flood periods were measured using the advanced 'MLPClassifeir' approach of Deep Learning, developed by GRASS GIS which employs the Pythons's algorithms of Scikit-Learn library of machine learning. The increase of flooded inundated areas in the wetlands of the Ganges, the development of mangrove forests and the fragmentations of the forest land cover types were detected on the images for scenes processed with DL modules of GRASS GIS. The tested images included multispectral bands, and were collected on the cloud-free dates. The notable changes in the flooded areas were shown by the year 2021 with properties of land cover types evaluated with image classification methods after seven flooding period which is impacted by the monsoon effects in southern Bangladesh.

%----------- subsection ----------->
\subsection{K-means classification outcomes}

Flooded areas in Ganges River Delta detected and mapped using MaxLike methods are presented as a series of the resulting maps in Figure \ref{fig08} and Figure \ref{fig09} and deep learning (DL) is presented in Figure \ref{fig10} and Figure \ref{fig11}. Maps of floods can be used as a broad base of environmental applications and flood hazard assessment. Their application can range from academia and research, industry and development, to governmental and social services for risk prevention. The findings presented in the study revealed a changing pattern of expanded floods in the post-flood period in Ganges Delta, and inequality in distribution of inundated areas covered by water by consequent years (2021, 2022 and 2023). 

The maps are generated using classified images for March showing the flooding event (in Figure \ref{fig08}) and Figure \ref{fig10}, as well as for November (Figure \ref{fig09} and Figure \ref{fig11}) showing the post-flood landscapes, respectively. Scripting models of GRASS GIS with open source available specifications catalyse further innovation in risk management and presents solutions for state-of-the-art mapping using programming approaches and machine learning modules. The following 10 land cover types were identified around the Ganges River Delta, Bangladesh: \emph{1); Water; 2) Wetlands and mudflats; 3) Mangrove forests; 4) Sandy areas; 5) Forests; 6) Croplands; 7) Grasslands; 8) Urban settlements; 9) Orchards; 10) Aquaculture.} \hl{Information on land cover types is derived from Rahman et al.} \cite{Rahman2020}\hl{, Sousa and Small} \cite{rs13234799} \hl{and Paszkowski et al.} \cite{Paszkowski2021}\hl{ and applied to the extent of the current study area.}

\begin{figure}[H]
\hspace{50pt}\makebox[0.90\linewidth][r]{%
	\begin{subfigure}[b]{.40\textwidth}
		\centering
		\includegraphics[width=0.98\textwidth]{Fig_08a.jpg}
		\caption{2021}
	\end{subfigure}%
	\begin{subfigure}[b]{.40\textwidth}
		\centering
		\includegraphics[width=0.98\textwidth]{Fig_08b.jpg}
		\caption{2022}
	\end{subfigure}%
	\begin{subfigure}[b]{.40\textwidth}
		\centering
		\includegraphics[width=0.98\textwidth]{Fig_08c.jpg}
		\caption{2023}
	\end{subfigure}%
}
\caption{Flooded landscapes in Ganges Delta, Bangladesh mapped using k-means clustering and MaxLike classification of GRASS GIS applied to the Landsat 8-9 OLI/TIRS images: \textbf{(a)}: March 2021; \textbf{(b)}: March 2022; \textbf{(c)}: March 2023.}\label{fig08}
\end{figure}

\begin{figure}[H]
\hspace{50pt}\makebox[0.90\linewidth][r]{%
	\begin{subfigure}[b]{.40\textwidth}
		\centering
			\includegraphics[width=0.98\textwidth]{Fig_09a.jpg}
		\caption{2021}
	\end{subfigure}%
	\begin{subfigure}[b]{.40\textwidth}
		\centering
		\includegraphics[width=0.98\textwidth]{Fig_09b.jpg}
		\caption{2022}
	\end{subfigure}%
	\begin{subfigure}[b]{.40\textwidth}
		\centering
		\includegraphics[width=0.98\textwidth]{Fig_09c.jpg}
		\caption{2023}
	\end{subfigure}%
}
\caption{Post-flood landscapes in Ganges Delta, Bangladesh mapped using k-means clustering and MaxLike classification of GRASS GIS applied to the Landsat 8-9 OLI/TIRS images: \textbf{(a)}: November 2021; \textbf{(b)}: November 2022; \textbf{(c)}: November 2023.}\label{fig09}
\end{figure}

%----------- subsection ----------->
\subsection{Deep Learning outcomes}

The supervised learning models of GRASS GIS based on the MLPClassifier demonstrated accurate mapping and refined visualisation functionality, Figure \ref{fig10} and Figure \ref{fig11}. While this method was developed to serve a specific goal of image processing, its purpose is applicable as solutions to environmental monitoring with outcomes of detected flooded areas as presented above. Hence, the research outcomes have provided insights into the changing landscapes of the post-flooded period of Ganges Delta. 

Time-series analysis of the Ganges Delta using Landsat 8-9 OLI/TIRS dataset enabled by the USGS was used to plots and graphically display the changes in the flood and post-flood periods for southern Bangladesh for three test periods: 2021, 2022 and 2023. The data for 2014 were used as training polygon for deep learning modelling which requires the classification seed. As demonstrated in this paper, GRASS GIS presents an open-source mapping product that shows high potential for its adoption in environmental monitoring and mapping through embedded modules of machine learning and deep learning that support accurate image classification through advanced computer vision algorithms of pattern recognition and image analysis. 

\begin{figure}[H]
\hspace{50pt}\makebox[0.90\linewidth][r]{%
	\begin{subfigure}[b]{.40\textwidth}
		\centering
			\includegraphics[width=0.98\textwidth]{Fig_10a.jpg}
		\caption{2021}
	\end{subfigure}%
	\begin{subfigure}[b]{.40\textwidth}
		\centering
		\includegraphics[width=0.98\textwidth]{Fig_10b.jpg}
		\caption{2022}
	\end{subfigure}%
	\begin{subfigure}[b]{.40\textwidth}
		\centering
		\includegraphics[width=0.98\textwidth]{Fig_10c.jpg}
		\caption{2023}
	\end{subfigure}%
}
\caption{Flooded landscapes in Ganges Delta, Bangladesh mapped using Deep Learning (DL) classification MLPC of GRASS GIS applied to the Landsat 8-9 OLI/TIRS images: \textbf{(a)}: March 2021; \textbf{(b)}: March 2022; \textbf{(c)}: March 2023.}\label{fig10}
\end{figure}

\begin{figure}[H]
\hspace{50pt}\makebox[0.90\linewidth][r]{%
	\begin{subfigure}[b]{.40\textwidth}
		\centering
			\includegraphics[width=0.98\textwidth]{Fig_11a.jpg}
		\caption{2021}
	\end{subfigure}%
	\begin{subfigure}[b]{.40\textwidth}
		\centering
		\includegraphics[width=0.98\textwidth]{Fig_11b.jpg}
		\caption{2022}
	\end{subfigure}%
	\begin{subfigure}[b]{.40\textwidth}
		\centering
		\includegraphics[width=0.98\textwidth]{Fig_11c.jpg}
		\caption{2023}
	\end{subfigure}%
}
\caption{Post-flood landscapes in Ganges Delta, Bangladesh mapped using Deep Learning (DL) classification MLPC of GRASS GIS applied to the Landsat 8-9 OLI/TIRS images: \textbf{(a)}: November 2021; \textbf{(b)}: November 2022; \textbf{(c)}: November 2023.}\label{fig11}
\end{figure}

\begin{table}[H]
  \centering
  \begin{tabularx}{\textwidth}{p{0.9cm}p{0.9cm}p{0.9cm}p{0.9cm}p{0.9cm}p{0.9cm}p{0.9cm}p{0.9cm}p{0.9cm}p{0.9cm}p{0.9cm}}
  \toprule
  \multirow{2}{*}{Year} & \multicolumn{10}{c}{Classes of land cover types in March: Ganges-Brahmaputra River Delta} \\
     & 1 & 2 & 3 & 4 & 5 & 6 & 7 & 8 & 9 & 10 \\
    \midrule
    \textbf{2021} & 717 & 252 & 434 & 758 & 1005 & 1035 & 758 & 744 & 927 & 219 \\
    \textbf{2022} &  749 & 236 & 489 & 803 & 1030 & 1100 & 557 & 744 & 800 & 336 \\
     \textbf{2023} &  707 & 343 & 1055 & 381 & 1032 & 997 & 702 & 649 & 796 & 180 \\
\bottomrule
\end{tabularx}
  \caption{Flood period: estimated classes of land cover types for March period for years 2021, 2022 and 2023 in Ganges River Delta}
  \label{tab02}
\end{table}

\begin{table}[H]
  \centering
  \begin{tabularx}{\textwidth}{p{0.9cm}p{0.9cm}p{0.9cm}p{0.9cm}p{0.9cm}p{0.9cm}p{0.9cm}p{0.9cm}p{0.9cm}p{0.9cm}p{0.9cm}}
  \toprule
  \multirow{2}{*}{Year} & \multicolumn{10}{c}{Classes of land cover types in November: Ganges-Brahmaputra River Delta} \\
     & 1 & 2 & 3 & 4 & 5 & 6 & 7 & 8 & 9 & 10 \\
    \midrule
    \textbf{2021} & 689 & 402 & 579 & 565 & 1228 & 1109 & 571 & 849 & 647 & 208 \\
    \textbf{2022} & 700 & 468 & 365 & 587 & 1066 & 1048 & 683 & 942 & 683 & 326 \\
     \textbf{2023} & 715 & 511 & 231 & 599 & 1351 & 1132 & 718 & 707 & 685 & 310 \\
\bottomrule
\end{tabularx}
  \caption{Post-flood period: estimated classes of land cover types for November period for years 2021, 2022 and 2023 in Ganges River Delta}
  \label{tab03}
\end{table}

The corrections of the presented framework and Python-based algorithm of GRASS GIS using machine learning (ML) approach which uses MLPClassifier of Deep Learning (DL) presents an advanced tool for automated recognition of land cover types on the satellite images and monitoring floods in the prone to floods regions of the southern Bangladesh. The percentage of the correctly classified images was calculated and summarised in the matrices showing classes separability (Matrix \ref{eq1}, Matrix \ref{eq2} and Matrix \ref{eq3} for years 2021, 2022 and 2023, respectively) and evaluated in the accuracy assessment. 

Class separability matrices of land cover classes identified over Ganges River Delta which are presented in the matrices below: Matrix \ref{eq1} shows class separability matrices for March and November 2021; Matrix \ref{eq2} shows class separability matrices for March and November 2022, and Matrix \ref{eq3} shows class separability matrices for March and November 2023. Class separability matrices evaluated the assignment of the pixels distributed to various land cover classes over the landscapes of Ganges River Delta, hence it serves as a function of the regularisation parameters used for distinguishability of pixels according to their spectral reflectances. In this way, class separability matrices perform an iterative search for cells within the raster matrix of the image considering a range of values of pixels assigned to different land cover types.

\begin{equation}\label{eq1}
\setlength\arraycolsep{3pt}
\footnotesize
\begin{vmatrix}
    & 1 & 2 & 3 & 4 & 5 & 6 & 7 & 8 & 9 & 10 \\
1 & 0 &  &  &  &  &  &  &  &  &  \\
2 & 1.3 & 0 &  &  &  &  &  &  &  &  \\
3 & 2.2 & 1.0 & 0 &  &  &  &  &  &  &  \\
4 & 3.0 & 1.5 & 0.7 & 0 &  &  &  &  &  &  \\
5 & 3.2 & 1.2 & 1.2 & 1.2 & 0 &  &  &  &  &  \\
6 & 3.4 & 1.9 & 1.3 & 0.7 & 1.1 & 0 &  &  &  &  \\  
7 & 3.6 & 2.4 & 1.7 & 1.2 & 1.6 & 0.7 & 0 &  &  &  \\  
8 & 3.7 & 1.7 & 1.8 & 1.7 & 0.7 & 1.4 & 1.7 & 0 &  &  \\
9 & 3.4 & 1.6 & 2.3 & 2.5 & 1.3 & 2.3 & 2.6 & 1.0 & 0 &  \\   
10 & 3.8 & 2.2 & 2.8 & 2.8 & 2.0 & 2.7 & 3.0 & 1.8 & 1.1 & 0 \\      
\end{vmatrix}
 %
 \hspace{10pt}
 \begin{vmatrix}
    & 1 & 2 & 3 & 4 & 5 & 6 & 7 & 8 & 9 & 10 \\
1 & 0 &  &  &  &  &  &  &  &  &  \\
2 & 1.9 & 0 &  &  &  &  &  &  &  &  \\
3 & 2.6 & 1.0 & 0 &  &  &  &  &  &  &  \\
4 & 2.9 & 2.2 & 1.4 & 0 &  &  &  &  &  &  \\
5 & 3.9 & 2.0 & 0.9 & 1.2 & 0 &  &  &  &  &  \\
6 & 4.6 & 2.4 & 1.4 & 1.5 & 0.6 & 0 &  &  &  &  \\  
7 & 4.2 & 2.8 & 1.9 & 1.3 & 1.3 & 1.0 & 0 &  &  &  \\  
8 & 4.1 & 2.9 & 1.8 & 0.8 & 1.3 & 1.3 & 0.7 & 0 &  &  \\
9 & 3.9 & 3.3 & 2.5 & 1.0 & 2.2 & 2.3 & 1.5 & 1.1 & 0 &  \\   
10 & 3.5 & 3.3 & 2.6 & 1.5 & 2.3 & 2.3 & 1.8 & 1.5 & 0.9 & 0 \\       
 \end{vmatrix} 
\end{equation}

\begin{equation}\label{eq2}
\setlength\arraycolsep{3pt}
\footnotesize
\begin{vmatrix}
   & 1 & 2 & 3 & 4 & 5 & 6 & 7 & 8 & 9 & 10 \\
1 & 0 &  &  &  &  &  &  &  &  &  \\
2 & 1.5 & 0 &  &  &  &  &  &  &  &  \\
3 & 2.4 & 1.0 & 0 &  &  &  &  &  &  &  \\
4 & 3.3 & 1.7 & 0.7 & 0 &  &  &  &  &  &  \\
5 & 3.4 & 1.3 & 1.3 & 1.3 & 0 &  &  &  &  &  \\
6 & 3.8 & 2.2 & 1.3 & 0.7 & 1.2 & 0 &  &  &  &  \\  
7 & 3.6 & 2.4 & 1.6 & 1.1 & 1.5 & 0.6 & 0 &  &  &  \\  
8 & 3.9 & 1.8 & 1.9 & 1.8 & 0.7 & 1.4 & 1.6 & 0 &  &  \\
9 & 3.9 & 1.7 & 2.5 & 2.7 & 1.4 & 2.4 & 2.5 & 1.1 & 0 &  \\   
10 & 3.9 & 2.1 & 2.8 & 2.9 & 2.0 & 2.7 & 2.8 & 1.7 & 1.0 & 0 \\      
\end{vmatrix}
 %
 \hspace{10pt}
 \begin{vmatrix}
    & 1 & 2 & 3 & 4 & 5 & 6 & 7 & 8 & 9 & 10 \\
1 & 0 &  &  &  &  &  &  &  &  &  \\
2 & 1.8 & 0 &  &  &  &  &  &  &  &  \\
3 & 2.4 & 1.8 & 0 &  &  &  &  &  &  &  \\
4 & 2.7 & 0.9 & 1.3 & 0 &  &  &  &  &  &  \\
5 & 3.8 & 1.8 & 1.4 & 0.8 & 0 &  &  &  &  &  \\
6 & 4.5 & 2.2 & 1.8 & 1.3 & 0.6 & 0 &  &  &  &  \\  
7 & 4.1 & 2.5 & 1.5 & 1.7 & 1.3 & 1.0 & 0 &  &  &  \\  
8 & 4.1 & 2.5 & 1.0 & 1.5 & 1.1 & 1.2 & 0.7 & 0 &  &  \\
9 & 3.9 & 2.9 & 0.8 & 2.3 & 2.2 & 2.3 & 1.4 & 1.2 & 0 &  \\   
10 & 3.0 & 2.8 & 1.2 & 2.3 & 2.3 & 2.4 & 1.8 & 1.7 & 0.8 & 0 \\       
 \end{vmatrix} 
\end{equation}

\begin{equation}\label{eq3}
\setlength\arraycolsep{3pt}
\footnotesize
\begin{vmatrix}
    & 1 & 2 & 3 & 4 & 5 & 6 & 7 & 8 & 9 & 10 \\
1 & 0 &  &  &  &  &  &  &  &  &  \\
2 & 1.4 & 0 &  &  &  &  &  &  &  &  \\
3 & 3.0 & 1.0 & 0 &  &  &  &  &  &  &  \\
4 & 3.0 & 1.3 & 1.4 & 0 &  &  &  &  &  &  \\
5 & 3.8 & 1.7 & 0.7 & 1.6 & 0 &  &  &  &  &  \\
6 & 4.1 & 1.7 & 1.2 & 0.9 & 1.1 & 0 &  &  &  &  \\  
7 & 3.9 & 2.1 & 1.3 & 2.1 & 0.7 & 1.5 & 0 &  &  &  \\  
8 & 4.4 & 2.2 & 1.7 & 1.2 & 1.4 & 0.7 & 1.4 & 0 &  &  \\
9 & 4.2 & 2.3 & 2.3 & 1.1 & 2.3 & 1.4 & 2.4 & 1.1 & 0 &  \\   
10 & 3.9 & 2.6 & 2.6 & 1.8 & 2.5 & 2.0 & 2.6 & 1.7 & 1.1 & 0 \\      
\end{vmatrix}
 %
 \hspace{10pt}
 \begin{vmatrix}
    & 1 & 2 & 3 & 4 & 5 & 6 & 7 & 8 & 9 & 10 \\
1 & 0 &  &  &  &  &  &  &  &  &  \\
2 & 1.8 & 0 &  &  &  &  &  &  &  &  \\
3 & 2.1 & 1.4 & 0 &  &  &  &  &  &  &  \\
4 & 3.1 & 1.1 & 1.2 & 0 &  &  &  &  &  &  \\
5 & 3.9 & 1.9 & 1.3 & 0.8 & 0 &  &  &  &  &  \\
6 & 4.2 & 2.3 & 1.7 & 1.2 & 0.6 & 0 &  &  &  &  \\  
7 & 3.5 & 2.5 & 0.9 & 1.8 & 1.4 & 1.5 & 0 &  &  &  \\  
8 & 4.1 & 2.6 & 1.7 & 1.8 & 1.3 & 0.9 & 1.2 & 0 &  &  \\
9 & 4.0 & 2.9 & 1.4 & 2.2 & 1.7 & 1.6 & 0.6 & 0.9 & 0 &  \\   
10 & 3.1 & 2.7 & 1.6 & 2.2 & 1.9 & 1.9 & 1.0 & 1.6 & 0.9 & 0 \\       
 \end{vmatrix} 
\end{equation}
 
%----------- section ----------->
\section{\textcolor{blue}{Discussion}}

\hl{Floods are a major threat to human communities all over the world, with serious socio-economic and environmental consequences, given that they are the leading cause of natural disaster losses. Related research shows that one in four people worldwide is considered to be at significant risk from flooding. Bangladesh, which occupies low-lying flood and tidal plains, has one of the world's largest and most disaster-prone deltas. In Bangladesh in particular, the absolute number of people at risk of flooding is 94.42 million, or 57.5\% of the country's population (the second highest in the world). }Floods in Ganges River Delta present \hl{one of }the most catastrophic problems at the national level in Bangladesh \cite{Szabo}. \hl{With this regards, }preventive mapping of flooded areas presents the societal importance for the social system \hl{of Bangladesh }and maintaining the sustainability of both natural landscapes and social infrastructure. 

\hl{Social awareness} of \hl{the} flood hazards \hl{can help mitigating the risk associated with floods. Hence, remote sensing (RS) data processed using advanced methods and visualised on} maps is an essential tool in addressing environmental challenges through \hl{mapping the} areas at risks, cartographic display of \hl{the }inundated coastal and prognosis of flood in deltaic areas based on \hl{the }time series analysis. \hl{To this end, the advanced} methods \hl{of deep learning (DL) used }for operative mapping of flooded areas \hl{support} publicly accessible, modifiable, and distributable open-source cartographic products. \hl{Such products can be used for preventive monitoring of the areas at risk, }prognosis of flood events \hl{and hydrologic modelling}. Hence, operative monitoring and forecasting of \hl{flood} hazards can effectively reduce the catastrophic consequences of floods in Ganges River Delta. 

Advanced methods of deep learning (DL) for accurate \hl{processing of }satellite images \hl{are} possible using open-source GRASS GIS \hl{which presents an} effective method to retrieve geo-information. Deep learning in GIS domain enables to increase the accuracy and automation of mapping through the algorithms of artificial intelligence employed in the cartographic workflow. 
Land cover changes in the Ganges River Delta were determined over the evaluated period using time series analysis of the satellite images Landsat, and compared for flood- and post-flood periods. The results indicated the increase in wetlands and inundated areas during the post-flood period and demonstrated the effectiveness of the DL approach integrating GRASS GIS and scripting cartographic approaches in tasks of environmental monitoring in tropical regions, such as coastal regions of Bangladesh. 

\section{\textcolor{blue}{Conclusions}}
To conclude, this study \hl{demonstrated} the capacity of sequential analysis of the remote sensing data to identify changes in the coastal and wetland landscapes of Ganges River Delta over the last three years (2021-2023) for comparison of flood event with post-flood landscapes, illustrates how the satellite images can help identify extent of floods, and affected areas, and shows what role deep learning techniques of the GRASS GIS can play in the future technology landscapes in similar studies on vulnerability assessment and mapping floods. Importantly, this paper presented a case of flood mapping in Ganges Delta using the integration of open source data and methods: the dataset was obtained from the publicly available repository of the USGS and processed using freely available GRASS GIS software. Such approach is essential for the environmentalists from the developing countries such as Bangladesh, because open source data and tools can motivate researchers to continue and expand research using the proposed techniques. In future directions, this works can be extended for further environmental monitoring of the Ganges Delta by producing maps on other periods using other EO datasets. 

\hl{This} research has demonstrated an example of using \hl{the advanced} approach \hl{of ML }in cartography and \hl{RS} data processing. Thus, it has led to a greater understanding of the importance of using advanced methods of machine learning for environmental monitoring. The application of the ANN as a branch of machine learning methods, presented here as a case of MLP classifier algorithm demonstrated that such technique can be successfully applied in similar studies which require RS data for spatio-temporal analysis. Thus, the displayed time series of maps on Ganges Delta showing floods in March and post-flooded landscapes in November highlighted the benefits achieved by remote sensing data visualization for analysis of flood extent and comparison of climatic effects by recent consecutive years: 2021, 2022 and 2023. The spectral properties of the multispectral satellite images Landsat OLI/TIRS  were evaluated with regards to the land cover change development in the inundated areas of the Ganges River Delta during the period of flood and the post-restoration period. Using ML methods of GRASS GIS remote sensing software demonstrated a fine perspectives for image processing and environmental data analysis. We evaluated the proposed ML framework of GRASS GIS on the Landsat scenes, however, other satellite images can also be applied in future similar studies.

%----------- section ----------->

\institutionalreview{Not applicable.}

\informedconsent{Not applicable.}

\dataavailability{Not applicable} 

\acknowledgments{The authors thank the reviewers for reading and review of this manuscript.}

\conflictsofinterest{The author declares that they have no known competing financial interests neither personal relationships that could have influenced the work reported in this article.} 

\abbreviations{Abbreviations}{
The following abbreviations are used in this manuscript:\\

\noindent 
\begin{tabular}{@{}ll}
ANN & Artificial Neural Networks \\
DL & Deep Learning \\
DN & Digital Number \\
EO & Earth Observation \\
GEBCO & General Bathymetric Chart of the Oceans \\
GMT & Generic Mapping Tools \\
GRASS & Geographic Resources Analysis Support System \\
GIS & Geographic Information System \\
Landsat OLI/TIRS & Landsat Operational Land Imager and Thermal Infrared Sensor \\
ML & Machine Learning \\
MLP & MultiLayer Perceptron \\
NIR & Near Infrared \\ 
RF & Random Forest \\
RS & Remote Sensing \\
SVM & Support Vector Machines \\
SWIR & Shortwave Infrared \\
USGS & United States Geological Survey \\
UTM & Universal Transverse Mercator \\
WGS84 & World Geodetic System 84\\
WRS & Worldwide Reference System  
\end{tabular}
}

\appendixtitles{no} % Leave argument "no" if all appendix headings stay EMPTY (then no dot is printed after "Appendix A"). If the appendix sections contain a heading then change the argument to "yes".
\appendixstart
\appendix
\section[\appendixname~\thesection]{Metadata for satellite images }

\begin{table}[H]
\footnotesize
\caption{Metadata for the Landsat 8-9 OLI/TIRS images obtained from the USGS} \label{tabApp}
\begin{adjustwidth}{-\extralength}{0cm}
    \begin{tabularx}{\fulllength}{llll}
    \toprule
       	\textbf{Data Set Attribute} & \textbf{Attribute Value} & \textbf{Attribute Value} & \textbf{Attribute Value} \\ \hline
        \textbf{Landsat Scene Identifier} & LC81370442023066LGN00 & LC81370442022079LGN00 & LC81370442021076LGN00 \\ \hline
        \textbf{Date Acquired} & \textbf{2023/03/07} & \textbf{2022/03/20} & \textbf{2021/03/17} \\ \hline
        \textbf{Roll Angle} & -1 & 0 & 0 \\ \hline
        \textbf{Start Time} & 2023-03-07 04:24:33 & 2022-03-20 04:24:26.058914 & 2021-03-17 04:24:23.120295 \\ \hline
        \textbf{Stop Time} & 2023-03-07 04:25:05 & 2022-03-20 04:24:57.828914 & 2021-03-17 04:24:54.890294 \\ \hline
        \textbf{Land Cloud Cover} & 3.12 & 0.01 & 0.03 \\ \hline
        \textbf{Scene Cloud Cover L1} & 2.97 & 0.01 & 0.03 \\ \hline
        \textbf{Ground Control Points Model} & 732 & 771 & 776 \\ \hline
        \textbf{Ground Control Points Version} & 5 & 5 & 5 \\ \hline
        \textbf{Geometric RMSE Model} & 5683 & 5615 & 5694 \\ \hline
        \textbf{Geometric RMSE Model X} & 3887 & 3857 & 3585 \\ \hline
        \textbf{Geometric RMSE Model Y} & 4146 & 4081 & 4424 \\ \hline
        \textbf{Processing Software Version} & LPGS\_16.2.0 & LPGS\_15.6.0 & LPGS\_15.4.0 \\ \hline
        \textbf{Sun Elevation L0RA} & 51.70012312 & 56.10505202 & 55.19368850 \\ \hline
        \textbf{Sun Azimuth L0RA} & 134.86314148 & 129.90704446 & 131.01649112 \\ \hline
        \textbf{TIRS SSM Model} & FINAL & FINAL & FINAL \\ \hline
        \textbf{Data Type L2} & OLI\_TIRS\_L2SP & OLI\_TIRS\_L2SP & OLI\_TIRS\_L2SP \\ \hline
        \textbf{Satellite} & 8 & 8 & 8 \\ \hline
        \textbf{Scene Center Lat DMS} & "23°06'46.58""N" & "23°06'45.29""N" & "23°06'46.01""N" \\ \hline
        \textbf{Scene Center Long DMS} & "90°23'29.83""E" & "90°23'14.57""E" & "90°24'53.42""E" \\ \hline
        \textbf{Corner Upper Left Lat DMS} & "24°09'00.29""N" & "24°09'00""N" & "24°09'02.38""N" \\ \hline
        \textbf{Corner Upper Left Long DMS} & "89°13'54.73""E" & "89°13'44.11""E" & "89°15'19.58""E" \\ \hline
        \textbf{Corner Upper Right Lat DMS} & "24°11'21.62""N" & "24°11'21.52""N" & "24°11'22.45""N" \\ \hline
        \textbf{Corner Upper Right Long DMS} & "91°30'23.18""E" & "91°30'12.56""E" & "91°31'48.22""E" \\ \hline
        \textbf{Corner Lower Left Lat DMS} & "22°01'28.45""N" & "22°01'28.20""N" & "22°01'30.32""N" \\ \hline
        \textbf{Corner Lower Left Long DMS} & "89°17'26.70""E" & "89°17'16.22""E" & "89°18'50.22""E" \\ \hline
        \textbf{Corner Lower Right Lat DMS} & "22°03'36.04""N" & "22°03'35.93""N" & "22°03'36.76""N" \\ \hline
        \textbf{Corner Lower Right Long DMS} & "91°31'47.39""E" & "91°31'36.95""E" & "91°33'11.12""E" \\ \hline
        \midrule
        \textbf{Data Set Attribute} & \textbf{Attribute Value} & \textbf{Attribute Value} & \textbf{Attribute Value} \\ \hline
        \textbf{Landsat Scene Identifier} & LC91370442023330LGN00 & LC91370442022327LGN01 & LC81370442021332LGN00 \\ \hline
        \textbf{Date Acquired} & \textbf{2023/11/26} & \textbf{2022/11/23} & \textbf{2021/11/28} \\ \hline
        \textbf{Roll Angle} & 0 & 0 & 0 \\ \hline
        \textbf{Start Time} & 2023-11-26 04:24:51 & 2022-11-23 04:25:01 & 2021-11-28 04:24:54.925489 \\ \hline
        \textbf{Stop Time} & 2023-11-26 04:25:23 & 2022-11-23 04:25:32 & 2021-11-28 04:25:26.695489 \\ \hline
        \textbf{Land Cloud Cover} & 0.02 & 0.20 & 0.40 \\ \hline
        \textbf{Scene Cloud Cover L1} & 0.02 & 0.19 & 0.38 \\ \hline
        \textbf{Ground Control Points Model} & 750 & 782 & 768 \\ \hline
        \textbf{Ground Control Points Version} & 5 & 5 & 5 \\ \hline
        \textbf{Geometric RMSE Model} & 6651 & 6490 & 6697 \\ \hline
        \textbf{Geometric RMSE Model X} & 4185 & 4191 & 4329 \\ \hline
        \textbf{Geometric RMSE Model Y} & 5170 & 4955 & 5110 \\ \hline
        \textbf{Processing Software Version} & LPGS\_16.3.1 & LPGS\_16.2.0 & LPGS\_15.5.0 \\ \hline
        \textbf{Sun Elevation L0RA} & 41.83496249 & 42.43660966 & 41.36291892 \\ \hline
        \textbf{Sun Azimuth L0RA} & 154.44654325 & 154.42300606 & 154.52745495 \\ \hline
        \textbf{TIRS SSM Model} & N/A & N/A & FINAL \\ \hline
        \textbf{Data Type L2} & OLI\_TIRS\_L2SP & OLI\_TIRS\_L2SP & OLI\_TIRS\_L2SP \\ \hline
        \textbf{Satellite} & 9 & 9 & 8 \\ \hline
        \textbf{Scene Center Lat DMS} & "23°06'45.65""N" & "23°06'46.48""N" & "23°06'45.22""N" \\ \hline
        \textbf{Scene Center Long DMS} & "90°22'20.75""E" & "90°23'11.94""E" & "90°24'08.21""E" \\ \hline
        \textbf{Corner Upper Left Lat DMS} & "24°08'48.98""N" & "24°08'50.28""N" & "24°09'01.33""N" \\ \hline
        \textbf{Corner Upper Left Long DMS} & "89°12'51.37""E" & "89°13'44.40""E" & "89°14'37.14""E" \\ \hline
        \textbf{Corner Upper Right Lat DMS} & "24°11'11.26""N" & "24°11'11.80""N" & "24°11'22.06""N" \\ \hline
        \textbf{Corner Upper Right Long DMS} & "91°29'19.50""E" & "91°30'12.67""E" & "91°31'05.70""E" \\ \hline
        \textbf{Corner Lower Left Lat DMS} & "22°01'27.01""N" & "22°01'28.20""N" & "22°01'19.67""N" \\ \hline
        \textbf{Corner Lower Left Long DMS} & "89°16'24.02""E" & "89°17'16.22""E" & "89°18'08.71""E" \\ \hline
        \textbf{Corner Lower Right Lat DMS} & "22°03'35.46""N" & "22°03'35.93""N" & "22°03'26.64""N" \\ \hline
        \textbf{Corner Lower Right Long DMS} & "91°30'44.60""E" & "91°31'36.95""E" & "91°32'29.36""E" \\ \hline
        \bottomrule
    \end{tabularx}
    \end{adjustwidth}
\end{table}

%%%%%%%%%%%%%%%%%%%%%%%%%%%%%%%%%%%%%%%%%%
\begin{adjustwidth}{-\extralength}{0cm}
%\printendnotes[custom] % Un-comment to print a list of endnotes

\reftitle{References}
%\nocite{*}

%=====================================
% References, variant A: external bibliography
%=====================================
\bibliography{Ganges.bib}

%%%%%%%%%%%%%%%%%%%%%%%%%%%%%%%%%%%%%%%%%%
\end{adjustwidth}
\end{document}
